\documentclass[12pt,a4paper,onecolumn]{article}

% Encoding and language
\usepackage[utf8]{inputenc}
\usepackage[T1]{fontenc}
\usepackage[english]{babel}

% Font
\usepackage{times}
\usepackage[normalem]{ulem}

% Set margins
\usepackage[a4paper, top=2cm, bottom=2cm, left=2.5cm, right=2.5cm]{geometry}

% Tabulation
\usepackage{tabto}

% Line spacing
\usepackage{setspace}
\onehalfspacing

% Page numbers
\usepackage{fancyhdr}
\pagestyle{fancy}
\fancyhf{} % Clear default header/footer
\rfoot{\thepage} % Right-aligned page numbers in footer

% Remove page number from cover/abstract/etc
\usepackage{titlesec}
\newcommand{\noPageNumber}{\thispagestyle{empty}}

% Prevent indentation
\setlength{\parindent}{0pt}
\setlength{\parskip}{0.5em}

% Figure management
\usepackage{graphicx}
\usepackage{caption}
\usepackage{subcaption}
\usepackage{wrapfig}

\usepackage{pgfplots}
\pgfplotsset{compat=1.16}
\def\mathdefault#1{#1}
\everymath=\expandafter{\the\everymath\displaystyle}

% Subfigure management
\usepackage{layouts}
%\usepackage{floatrow}
%\usepackage[label font=bf,labelformat=simple]{subfig}% <-- changed
%\floatsetup[figure]{style=plain,subcapbesideposition=top}


% Optional: headings formatting
\titleformat{\section}{\large\bfseries}{\thesection}{1em}{}
\titleformat{\subsection}{\normalsize\bfseries}{\thesubsection}{1em}{}

% Optional: remove widow/orphan lines
\widowpenalty=10000
\clubpenalty=10000

% Mathematics
\usepackage{amsmath,mathtools,amssymb} 

\begin{document}
 %% textwidth in cm: \printinunitsof{in}\prntlen{\textwidth}

% Cover page (no page number)
\noPageNumber
\begin{titlepage}
\centering
\vspace*{3cm}
{\Huge Evaluating the impact of maturation time variability in stage-structured host-parasitoid models (cette page de garde va changer) \par}

\begin{figure}[]
\centering
\includegraphics[width = 1\textwidth]{/home/mathisdr/Desktop/ParisTexas5.png}
\end{figure}

\vspace{2cm}
{\Large Mathis Drieu \par}
\vfill
February-June 2025
\end{titlepage}

% Abstract (no page number)
\noPageNumber
\begin{abstract}
xxx
\\
xxx
\\
\\
xxx
\\ 

\end{abstract}

% Table of Contents, List of Abbreviations (no page numbers)
\noPageNumber
\tableofcontents


% \listoffigures % If needed
% \listoftables  % If needed

% List of abbreviations
\section*{List of Abbreviations}
\noPageNumber
% Your content

% Main content starts here (with page numbers from 1)
\newpage
\setcounter{page}{1}



















\section{Funding : en attente d'un meilleur template?}
The internship is funded by the Cacolac project (OFB, Ecophyto...), which objective is to integratively develop agroecological methods of pest control. It aims to better understand the complementarity or antagonism between different operational levers used to control bio-aggressors in perennial crops and to optimize their combination according to the local context and landscape.


% References (excluded from page count)
\newpage
\noPageNumber
\begin{thebibliography}{99}
\bibitem{example} Author, \emph{Title}, Journal, Year.
\end{thebibliography}





% Appendix (excluded from page count)
\newpage
\noPageNumber
\appendix
\section{Appendix}

\subsection{Murdoch 1987 system}
\subsubsection{System construction : application of the Linear Chain Trick with a non-autonomous mortality term}
Note that Hurtado et al. 2018 as well as Lou and Sun 2022 give detailed justifications about the construction of continuous stage-structured models with an arbitrary maturation distribution and their convergence to a coupled system of ODEs in the case of the Erlang distribution. However, they do not include a non-autonomous mortality term. Let us thereby justify thoroughly how the studied LCT model was built analogously to the DDE model given by Murdoch (1987).

Let us consider the general form of the models used both in the ODE and DDE framework :  
\begin{equation}
\dot{H_{1}}(t) = B(t) - M_n(t) - \delta(t, ...) H_1(t)
\end{equation}

$H_1(t)$ is the density of juvenile hosts. $B(t)$ is the recruitment term of the juveniles, ie the birth rate. $M_n(t)$ is the maturation term. In the ODE framework, it would be of the form $M_{ODE}(t)= m_H H_1(t)$. In the DDE framework, it would be of the form $M_{DDE}(t) = B(t-\tau_H)S_{DDE}(t)$ with $S_{DDE}(t) := \exp\left( -\int_{t-\tau_H}^{t} \delta(u, ...) \, du \right)$ (Gurney and Nisbet 1983). 

$\delta_H(t, ...)$ is a mortality term. \textbf{We assume that $\delta$ can depend on both the densities of host and parasitoids}. Though, for clarity, $\delta(t, ...)$ will now be written in the form $\delta(t)$.

We give to the equation the initial conditions $H_1(0) = H_1^0$ and $B(-t) = B^{-}(t)$, $\delta(-t) = \delta^{-}(t)$, for $t \geq 0$.

To account for a arbitrary residence time distribution, the maturation term can be written in an integral form :
\begin{equation}
M_n(t) = \int_0^{+\infty} B(t-s)\varphi_n(s) \exp\left( -\int_{t-s}^{t} \delta(u) \, du \right) ds
\end{equation}
where $\varphi_n$ is the probability density function (PDF) of a gamma distribution of rate $\theta_H$ and shape n. Let's recall that

\begin{align}
\varphi_1(0) &= \theta_H, & \varphi_1' &= -\theta_H \varphi_1 \\
\varphi_n(0) &= 0, & \varphi_n' &= \theta_H (\varphi_{n-1} - \varphi_n) \quad \text{for } n \geq 2
\end{align}

Let us write :
\[
S(s, t) := \exp\left( -\int_{t-s}^{t} \delta(u) \, du \right)
\]
which satisfies :
\begin{equation}
S(0, t) = 1, \quad \partial_s S(s,t) + \partial_t S(s,t) = -\delta(t) S(s,t)
\end{equation}

We can thus write :
\begin{equation}
M_n(t) = \int_0^{+\infty} B(t-s)\varphi_n(s)S(s,t) \, ds
\end{equation}

By derivation :
\[
M_n'(t) = \int_0^{+\infty} B'(t-s)\varphi_n(s)S(s,t) \, ds + \int_0^{+\infty} B(t-s)\varphi_n(s) \partial_t S(s,t) \, ds
\]

The first term, with an integration by part in s and using $S(0,t) = 1$, gives :
\begin{align*}
\int_0^{+\infty} B'(t-s)\varphi_n(s)S(s,t) \, ds &= -[B(t-s) \varphi_n(s) S(s,t) ]_{0}^{\infty} + \int_0^{+\infty} B(t-s)\frac{d}{ds}(\varphi_n(s)S(s,t)) \, ds \\
&\quad = -0 + \varphi_n(0) B(t) + \int_0^{+\infty} B(t-s)\varphi_n'(s) S(s,t) \, ds \\
&\quad + \int_0^{+\infty} B(t-s)\varphi_n(s) \partial_s S(s,t) \, ds
\end{align*}

With $(4)$, we have :
\begin{align*}
&\int_0^{+\infty} B(t-s)\varphi_n(s) \partial_s S(s,t) \, ds  + \int_0^{+\infty} B(t-s)\varphi_n(s) \partial_t S(s,t) \, ds \\
&\quad = - \delta(t) \int_0^{+\infty} B(t-s)\varphi_n(s) S(s,t)  \, ds \\
&\quad = - \delta(t) M_n(t)
\end{align*}

we finally get :
\[
M_n'(t) = \varphi_n(0) B(t) + \int_0^{+\infty} B(t-s)\varphi_n'(s) S(s,t) \, ds - \delta(t) M_n(t)
\]

In the case $n = 1$, i.e in the ODE framework, using $(3)$ : 
\begin{align*}
M_{1}'(t) &= M_{ODE}'(t)= -\theta_H B(t) + \int_0^{+\infty} B(t-s)(-\theta_H \varphi_1(s)) S(s,t) \, ds - \delta(t) M_1(t) \\
&= \theta_H B(t) - (\theta_H + \delta(t)) M_1(t)
\end{align*}

And for $n \geq 2$, using $(4)$  :
\[
M_n'(t) = \theta_H M_{n-1}(t) - (\theta_H + \delta(t)) M_n(t)
\]

Thus, \(M_n\) is solution of the system :
\begin{equation}
\begin{cases}
\dot{M}_1(t) = \theta_H B(t) - (\theta_H + \delta(t)) M_1(t) \\
\dot{M}_k(t) = \theta_H M_{k-1}(t) - (\theta_H + \delta(t)) M_k(t), \quad k \in \{2, \dots, n\}
\end{cases}
\end{equation}

The initial conditions of this system  $(M_1(0), \dots, M_N(0))$ are given by $(2)$, with $t = 0$.
Now, let us write, for \(k \in \{1, \dots, N\}\), $H_{1,k}(t) := \frac{M_k(t)}{\theta_H}$. Then : 
\begin{equation}
\begin{cases}
\dot{H_{1,1}}(t) = B(t) - (\theta_H + \delta(t)) H_{1,1}(t) \\
\dot{H_{1,1}}(t) = \theta_H H_{1, k-1}(t) - (\theta_H + \delta(t)) H_{1,k}(t), \quad k \in \{2, \dots, n\}
\end{cases}
\end{equation}

By summing all the lines, and writing $y(t) := H_{1,1}(t) + \dots + H_{1,n}(t)$, we get :
\[
\dot{y}(t) = B(t) - \theta_H H_{1,n}(t) - \delta(t) y(t) = B(t) - M_n(t) - \delta(t) y(t)
\]

Thus, $H_{1,1} + \dots + H_{1,n}$ does satisfy the initial condition $(1)$.

In particular, if $H_{1,1}(0) + \dots + H_{1,n}(0) =  H_1^0$, then $H_1(t) = y(t) = H_{1,1}(t) + \dots + H_{1,n}(t)$ and $(1)$ is equivalent to $(8)$. This is called the Linear Chain Trick.

Applying the same method on the parasitoid maturation term with Erlang rate $\theta_P$ and shape $m$, choosing a linear parasitism term of parasitism assuming a random and homogeneous search of hosts by the parasitoids (May 1978), and assuming a linear birth rate, the final LCT system can be rewritten as a system of coupled ODEs:

\begin{equation}
\begin{dcases}
\dot{H_{1,1}}(t) = H_{1,1}(t)(-a P_2(t) - \delta_{H1} - \theta_H) + r H_2(t) \\
\vdots \\
\dot{H_{1,i}}(t) = \theta_H(H_{1,i-1}(t) - H_{1,i}(t)) - \delta_{H1} H_{1,i}(t) - aP_2H_{1,i}(t), \quad i \in (2, 3, \dots, n) \\
\vdots \\
\dot{H_{2}}(t) = \theta_H H_{1,n}(t) - \delta_{H2} H_{2}(t) \\
\dot{P_{1,1}}(t) = P_{1,1}(t)(- \delta_{P1} - \theta_P) + a P_2(t) \sum_{i=1}^{n} H_{1,i}(t) \\
\vdots \\
\dot{P_{1,i}}(t) = \theta_P(P_{1,i-1}(t) - P_{1,i}(t)) - \delta_{P1} P_{1,i}(t), \quad i \in (2, 3, \dots, m) \\
\vdots \\
\dot{P_{2}}(t) = \theta_P P_{1,m}(t) - \delta_{P2} P_{2}(t) \\
\end{dcases}
\end{equation}

\subsubsection{Equilibria}

Setting the right side of the model system equal to zero and solving for equilibrium yields the following non-trivial equilibrium densities :
\begin{align}
H_1^* &= \frac{\delta_{P2}}{a(f_P)^m} \nonumber \\
H_2^* &= H_1^* * \frac{\theta_{H} (\rho^{\frac{1}{n}}-1)}{\delta_{H2} (\rho-1)} \nonumber \\
P_1^* &= P_2^* * \frac{\alpha_P \delta_{P2}}{\theta_P (f_P)^{m-1}}  \\
P_2^*&= \frac{\theta_H}{a f_H} (\rho^{\frac{1}{n}}f_H -1) \nonumber
\end{align}

With :
\begin{align}
f_* &:= \frac{\theta_*}{\theta_* + \delta_*}, \quad \text{$_*$ equals H,P} \nonumber  \\
s_H^* &:= \frac{\theta_H}{\theta_H + \delta_H + aP_2^*} = \rho^{-\frac{1}{n}} \nonumber \\
\alpha_H^* &:= \sum_{i=0}^{n-1} (s_H^*)^i   \\
\alpha_P^* &:= \sum_{i=0}^{m-1} (f_P^*)^i \nonumber
\end{align}

These equilibria are strictly positive if :
\begin{equation}
(\frac{\beta}{\delta_{H2}})^{\frac{1}{n}} f_H  > 1
\end{equation}

\subsubsection{Parametrisation}
To find the mean development time, or expected lifespan of juveniles that reach maturity, the Erlang distribution properties can be used.
For clarity, let us reason for hosts in the following equations. Parasitoid terms are set in the same way.  Let us recall that the maturation time is Erlang-distributed with rate $\theta_H$ and shape n. For juveniles, the mean maturation time $\hat{\tau_{H}}$, the variance of maturation time $\sigma^2_{H}$ and the coefficient of variation of maturation time $CV_{H}$ would be : 
\begin{align}
\hat{\tau_{H}} = \frac{n}{\theta_{H}} \nonumber \\
\sigma^2_{H} = \frac{n}{(\theta_{H})^2} \\
CV_{H} = \frac{1}{\sqrt{n}} \nonumber 
\end{align}

Setting  $\hat{\tau_{H}} = 1/m_H$, the maturation term becomes :
\begin{equation}
\theta_H = n*m_H
\end{equation}

Also, setting $\rho$ the mean number of offspring yielded by a host individual during its adult lifetime, we get :
\begin{equation}
r = \rho \delta_{H2}
\end{equation}

However, in the previous lines, juvenile intrinsic mortality was not taken in account. Alternatively, a continuous-time stochastic agent-based model can be built. This intuitively accounts for the stochastic process lying behind the model assumptions.

Consider a juvenile agent that either matures or dies. Maturation time $\tau_m$ is drawn from the Erlang distribution $\varphi(\theta_H, n)$ and death time $\tau_d$ is drawn from the Exponential distribution $\mathcal{E}(\delta_H)$. If $\tau_m \leq \tau_d$ then the agent matures. Once adult, new juveniles are produced from a Poisson point process of rate $\beta$.

Thus, let $f_{surv,H1}$ be the mean proportion of surviving host individuals and (13) become :

\begin{align}
\hat{\tau_{H}} = \frac{n}{\theta_{H} + \delta_{H}} \nonumber \\
\sigma^2_{H1} = \frac{n}{(\theta_{H} + \delta_{H})^2} \\
CV_{H} = \frac{1}{\sqrt{n}} \\
f_{surv,H1} = (\frac{\theta_{H}}{\theta_{H} + \delta_{H}})^n
\end{align}

Still setting $\hat{\tau_{H}} = 1/m_H$, the maturation and intrinsic mortality terms then become :
\begin{equation}
\begin{alignedat}{2}
\theta_{H} = n * m_{H}* (f_{surv,H1})^{1/n} \\
\delta_{H} = n* m_{H}* (1-(f_{surv,H1})^{1/n})
\end{alignedat}
\end{equation}

We notice that the existence of a positive non-trivial equilibrium can be rewritten as :
$$\rho^{\frac{1}{n}}f_H = (\rho f_{surv,H1})^{\frac{1}{n}} > 1 \qquad \impliedby \qquad \rho f_{surv,H1} > 1$$

\subsubsection{Convergence to ODE and DDE models}
If $n,m = (1,1)$, the Erlang distribution is equal to an exponential distribution and the system is equivalent to its simple ODE system counterpart. Then, $f_*$, $s_H$, $\alpha_*$ are all equal to 1 and equilibrium densities are those found in the simple ODE system. \\ \\

If $n,m\to\infty$, the Erlang distribution of maturation time converges to a dirac distribution and the system is equivalent to its DDE counterpart. Then, equibrium densities become : 
\begin{align}
H_1^* & \xrightarrow[n,m\to\infty]{} \frac{\delta_{P2}}{a\gamma} \nonumber \\
H_2^* & \xrightarrow[n,m\to\infty]{} \frac{ln(\rho)}{\rho-1} \frac{\delta_{P2} m_H}{a \gamma \delta_{H2}} \nonumber \\
P_1^* & \xrightarrow[n,m\to\infty]{} P_2^{* \infty} \frac{\gamma -1}{\gamma} \frac{\delta_{P2}}{m_P ln(\gamma)} = \frac{(\gamma-1)m_H(ln(\rho) - \frac{\delta_{H1}}{m_H})}{\gamma  a  m_P   ln(\gamma)} \\
P_2^*&  \xrightarrow[n,m\to\infty]{} \frac{m_H}{a}(ln(\rho) - \frac{\delta_{H1}}{m_H}) \nonumber  \\
\text{With} \quad \gamma & := e^{-\frac{\delta_{P1}}{m_P}} \sim_{m \to \infty} (f_P)^m \nonumber
\end{align}


Expectedly, they are strictly equivalent to equilibrium densities given in Murdoch 1987.

\subsection{Godfray 1989 system}

Setting the right side of the model system equal to zero and solving for equilibrium yields the following non-trivial equilibrium densities:
 \begin{align}
  H_1^* &= H_2^* * \frac{\theta_{H2}+\delta_{H2}+\Pi^*}{\theta_{H1}} \frac{\alpha_{H1}}{\alpha_{H2}} (f_{H1})^{1-n_1} \nonumber \\
  H_2^* &= P_2^* * \frac{\delta_{P2}}{(f_{P1})^{m} \Pi^*} \nonumber \\
  H_3^* &= H_4^* * \frac{\delta_{H4} \alpha_{H3}}{\theta_{H3}(f_{H3})^{n_3-1}}  \nonumber \\
  H_4^* &= H_1^* * \frac{\theta_{H1}+\delta_{H1}}{r \alpha_{H1}}  \\
  P_1^* &= H_2^* * \frac{\Pi^*}{\delta_{P1}}(1-(f_{P1})^{m}) \nonumber \\
  P_2^* &= \frac{k}{a}exp(\frac{\Pi*}{k}-1) \nonumber 
\end{align}

With :
 \begin{align}
  f_X^* &:= \frac{\theta_X}{\theta_X + \delta_X}, \quad \text{$X$ equals H1,H3,P1} \nonumber \\
  s_{H2}^* &:= \frac{\theta_{H2}}{\theta_{H2} + \delta_{H2} + \Pi^*} = [\frac{\delta_{H4}}{r}(f_{H1})^{-n1}(f_{H3})^{-n_3}]^{\frac{1}{n_2}}\nonumber \\
  \alpha_X^* &:= \sum_{i=0}^{n_X-1} (f_X^*)^i,  \quad \text{$X$ equals H1, H3,P1} \\
  \alpha_{H2}^* &:= \sum_{i=0}^{n_2-1} (s_{H2}^*)^i \nonumber \\
  \Pi^* &:= k*ln(1+\frac{aP_2^*}{k}) = \frac{\theta_{H2}}{f_{H2}} ((\frac{r}{\delta_{H4}}(f_{H1})^{n1}(f_{H3})^{n3})^{\frac{1}{n2}}f_{H2} -1)\nonumber 
 \end{align}


\subsection{Supplementary figures}
\renewcommand{\figurename}{SFig}
\setcounter{figure}{0}


\begin{figure}[!ht]
\centering	
\includegraphics[width=0.8\textwidth]{/home/mathisdr/Desktop/2024-2025/Stage/Internship_2025_LCT/model_host_parasitoid/figs/Murdoch_bifurcation_diagram_Other-par.pdf}
\caption{Bifurcation diagram of the system with parameters inspired from those in Murdoch 1987 : $\delta_{H1} = 0, \delta_{P1} = 1, \delta_{P2} = 8, \rho = 33, m_P = 2,  T_{H2} = 1/\delta_{H2} = 2, T_{H1} = 1/m_H = 1$.}
\label{fig:bifurc_suppl}
\end{figure}

\setcounter{figure}{-1}
\begin{figure}[!ht]
\centering	
\includegraphics[width=0.8\textwidth]{/home/mathisdr/Desktop/2024-2025/Stage/Internship_2025_LCT/model_host_parasitoid/figs/pbmtq_figure.png}
\caption{Exploring the effect of developmental (time) variability on host-parasitoid models using the linear chain trick.}
\label{fig:intro}
\end{figure}


\end{document}
